%\documentclass[runningheads]{llncs}
%\documentclass[orivec]{llncs}
\documentclass[runningheads, orivec]{llncs}
\usepackage[T1]{fontenc} 
\usepackage{graphicx}
\usepackage{nccmath}
\usepackage{marvosym} 
\usepackage{color}
\usepackage{hyperref}
\usepackage{algorithm}
\usepackage{algorithmic}
\usepackage[utf8]{inputenc}
\usepackage{csquotes}            
\usepackage{lipsum}                 
\usepackage{adjustbox}
\usepackage{siunitx}
\usepackage{float}
\usepackage{listings}
\usepackage{graphbox}
\usepackage{multirow}
\usepackage{booktabs} 
\usepackage{verbatim}
\usepackage[T1]{fontenc}
\usepackage[english]{babel}
\begin{document}
\title{A Simple and Effective Baseline for Attentional Generative Adversarial Networks}
% If the paper title is too long for the running head, you can set
% an abbreviated paper title here
    
\titlerunning{A Simple and Effective Baseline for Attentional GAN}

\author{
Mingyu Jin\inst{1*}\and 
Chong Zhang\inst{3*}\and
Qinkai Yu\inst{2} \and
Haochen Xue\inst{3} \and
Xiaobo Jin\inst{3}\textsuperscript{\Letter} \and
Xi Yang\inst{3}\textsuperscript{\Letter}
}
\vspace{-10pt}
\authorrunning{C. Zhang, M. Jin, Q. Yu, H. Xue, et al.}
\institute{Electrical and Computer Engineering Northwestern University, Evanston, Illinois, U.S.A.
\and
Department of Mathematics Science, University of Liverpool, Liverpool, U.K.
\and 
School of Advanced Technology, Xi‘an Jiaotong-Liverpool University, Suzhou, China 
\\
\email{u9o2n2@u.northwestern.edu} \quad \email{sgqyu9@liverpool.ac.uk}  \\
\email{\{Chong.zhang19, Haochen.Xue20\}@student.xjtlu.edu.cn} \\
\email{\{Xiaobo.Jin, Xi.Yang\}@xjtlu.edu.cn} 
\footnote{* Equal contribution. \\  \textsuperscript{\Letter} Corresponding author.}
}

\maketitle  % typeset the header of the contribution
\vspace{-15pt}

\begin{abstract}
Synthesising a text-to-image model of high-quality images by guiding the generative model through the Text description is an innovative and challenging task. In recent years, AttnGAN based on the Attention mechanism to guide GAN training has been proposed, SD-GAN, which adopts a self-distillation technique to improve the performance of the generator and the quality of image generation, and Stack-GAN++, which gradually improves the details and quality of the image by stacking multiple generators and discriminators. However, this series of improvements to GAN all have redundancy to a certain extent, which affects the generation performance and complexity to a certain extent. We use the popular simple and effective idea (1) to remove redundancy structure and improve the backbone network of AttnGAN. (2) to integrate and reconstruct multiple losses of DAMSM. Our improvements have significantly improved the model size and training efficiency while ensuring that the model's performance is unchanged and finally proposed our \textbf{SEAttnGAN}. Code is avalilable at \href{https://github.com/jmyissb/SEAttnGAN}{https://github.com/jmyissb/SEAttnGAN}.

\keywords{Text-to-image \and AttnGAN \and Simple and effective \and redundancy \and SEAttnGAN}
\end{abstract}

\vspace{-20pt}
\section{Introduction}

Text-to-image is a technique that converts natural language text into corresponding images, this technology covers some aspects of computer vision and natural language processing, and its application fields are very extensive, such as artistic creation, virtual scene generation, advertising design, game development, etc \cite{isola2017image} \cite{gatys2015neural} \cite{karras2017progressive} \cite{frolov2021adversarial} \cite{hald2020procedural}.  

\par Text-to-image is one of the most important applications of GAN (generative adversarial network) \cite{goodfellow2020generative}. In text-to-image synthesis, GAN’s generator receives text descriptions as input and generates images that match the descriptions. The discriminator is responsible for evaluating the authenticity of the generated image, and its goal is to distinguish the generated image from the real image. Generators and discriminators are trained through adversarial learning to gradually improve the quality of the generated images. The goal of GAN is to generate high-quality images of relevant content through semantics. This tasks presents two significant challenges: (1) Due to the instability of GAN structure, the generated images may become blurry and unrealistic. (2) Text information cannot be fully utilized by the generator, which limits the effectiveness of the text image fusion process and makes it difficult for the model to scale to higher resolution image synthesis. 

\par In order to solve the problems of traditional GAN, AttenGAN introduce the attention mechanism and deep attention multimodal similarity model in the generation work, this allows it to capture word level fine-grained information in sentences, allowing it to synthesize fine-grained details in different sub regions of the image by focusing on relevant words in natural language descriptions. However, the model structure of AttenGAN is complex and has a large number of parameters, this may lead to issues with high computational complexity. To solve the problem

\begin{itemize}
  \item We propse our optimized simplified version, called SEAttnGAN, our work improves the structure of AttnGAN and optimises the number of its parameters, and its performance maintains the same level as before optimisation. 
  \item Additionally, in order to achieve the corresponding effect even after model simplification and parameter reduction, we introduced semantics in the model. Our model can use the comparative loss of semantics as auxiliary information can guide the GAN generator to more accurately generate images that match the text description. 
  \item In order to solve the large and complex problem of the AttenGAN model, our work proposes a simplified AttenGAN that performs similarly to the original AttenGAN with significantly reduced structure and reduced parameter count. 
\end{itemize}

\par The overall paper is organized as follows: in Section 2, we introduce previous related work on text-to-image task. Section 3 gives a detailed introduction to our method; in Section 4, our algorithm will be compared with other algorithms. Subsequently, we summarize our algorithm and possible future research directions.
\vspace{-10pt}

\section{Related Work} 

Generative Adversarial Networks (GAN) were introduced in a seminal paper by Ian Goodfellow and his colleagues in 2014 \cite{goodfellow2020generative}. Goodfellow and his colleagues introduced GAN as a new approach to generative modelling. They showed that GAN could generate realistic images and that the generated images were difficult to distinguish from real images. This was a major breakthrough in generative modelling, as previous methods were often limited in their ability to generate high-quality data. In 2017, a paper in ICCV presented a method for generating lifelike images using only a written description in natural language \cite{dong2017semantic}. They change the structure of GAN and find it performs better than the method proposed by Reed et al. \cite{reed2016generative}. Specifically, they trained a style encoder network to invert the generator network.

While current methods can generate images with new attributes, they do not maintain the text-irrelevant details of the original image. In 2018 to address this issue, some authors at Yonsei University in NIPS 2018 introduced the Text-Adaptive Generative Adversarial Network (TAGAN), which generates semantically manipulated images while preserving the text-irrelevant details \cite{nam2018text}. This research paper proposes a solution to the problem of manipulating images based on natural language descriptions. The key idea of PPGN is to feed a random noise vector to the generative model and combine it with optimization techniques to guide the generative process \cite{nguyen2017plug}. 

LAPGAN is a generative model framework that combines Laplacian pyramid decomposition and GAN and aims to generate high-resolution images in a progressive manner \cite{denton2015deep}.



The network takes a multi-scale approach and generates images in a coarse-to-fine fashion, with each level of the model generating a low-resolution version of the image. The goal of MirrorGAN is to achieve text-to-image conversion by learning fine-grained correspondences \cite{qiao2019mirrorgan}. The proposed method promotes consistency between text and image through backward information propagation during generation, and provides additional supervision signals through the discriminator in the reverse generation process.

There have been several advancements in text-to-image generation using GAN. Here are some notable developments: Conditional GAN: In text-to-image synthesis, conditional GAN have been widely used where the generator takes a textual description as input and generates an image that matches the description \cite{mirza2014conditional}. Recent research has shown that conditioning the GAN on different factors such as object attributes, style, and layout can improve the quality and diversity of the generated images. The AttnGAN is a symbol of attention-based mechanisms \cite{xu2018attngan}.Attention-based mechanisms have been used to improve the quality and relevance of the generated images by focusing on the important parts of the text description. These mechanisms have been used to guide the generator to pay attention to the relevant words and phrases and use them to create the corresponding image. Multi-modal GANs have been developed to generate multiple images corresponding to a single text description, each depicting a different description aspect \cite{amirian2019social}. These models have been shown to generate diverse and high-quality images, capturing different interpretations of the same text description.

VAE(Variational Autoencoder) is a type of generative model that has been explored for text-to-image synthesis \cite{kingma2013auto}. In VAE-based text-to-image synthesis, the model learns a latent space representation of the input text descriptions, which can be used to generate images. One of the primary advantages of using VAE-based models is their interpretability. But, there is still room for improvement to generate photo-realistic images that match the input textual description. PixelCNN was a generative model generates images pixel by pixel, where each pixel is conditioned on the previous pixels \cite{van2016conditional}. The key idea behind PixelCNN is to use a convolutional neural network (CNN) to model the distribution of the image pixels. PixelCNN has been shown to be effective at generating high-quality images, it has also been used for text-to-image synthesis, where the image is generated conditioned on an input text description.

GAN-INT-CLS is the first application study that adds textual descriptions (i.e., sentence embedding vectors) as supervised signals to image generation \cite{reed2016generative}. The main goal of GAWWN is to estimate the scene layout of an image while generating it \cite{reed2016learning}. GAWWN improves the text-to-image performance by introducing Scene Layout Estimator in the image generation process \cite{zhang2017stackgan}. StackGAN (Stacked Generative Adversarial Networks) is used to gradually generate high-resolution images. This paper constructs a multi-stage generative model, which aims to solve the challenge of traditional GAN models in generating detail-rich images. Stackgan-v2 This paper improves and extends StackGAN by adding multiple generator and discriminator networks \cite{zhang2018stackgan++}, introducing a super-resolution network, and introducing a triangular loss function and a classification loss function. SD-GAN mainly solves the confounding problem of the traditional GAN model in representing factors in the latent space and introduces the mechanism of self-disentanglement learning to train the generator and discriminator networks \cite{yin2019semantics}. Training of the process-aided generator for "self-distillation" by augmenting the Distiller. SD-GAN mainly solves the confounding problem of the traditional GAN model in representing factors in the latent space, and introduces the mechanism of self-disentanglement learning to train the generator network and the discriminator network \cite{zhu2019dm}. Training of the process-aided generator for "self-distillation" by augmenting Distiller. 
\vspace{-10pt}
\section{Methodology}

In this paper, we simplified the network structure of AttnGAN. Here, suppose $z$ is a noise vector sampled from a standard normal distribution. $\bar{e}$  is a global sentence vector, and matrix of word vectors is $e$. For the input of the generator, we mix sentence features with noise. After data augmentation, the input is first reshaped through a fully connected layer. Then apply a series of Up Block layers to extract image features. The Up-Block layer is composed of an upsample layer, a redisual block, and DF-Blocks to fuse the text and image feature during the generation process. The image features obtained are used as  input and enter the attention layer together with the text features. Finally, the image features are re-extracted by a layer of Up-Block and converted to image features by a convolution layer. The detail architecrure shown in Figure \ref{method}\\
\begin{figure}[htbp]
  %\begin{center}
    \includegraphics[width=1.0\textwidth]{method.png}
  \caption{The architecture of the SEAttnGAN. FC: fully connected layer. UPBlock: residual block + DFBlock. DownBlock: downsample + residual block}
  \label{method}
  %\end{center}
\vspace{-20pt}
\end{figure}
We simplified the existing DFblock structure and used a new DFblock structure consisting of a two-layer linear mlp structure and an additional hidden layer (with ReLU activation) to further extract sentence features, map sentence features from 256 dimensional space to 64 dimensional space, and obtain new representation features. The new structure has the characteristic of lightweight, reducing the number of model parameters, simplifying the calculation process, and accelerating training speed. The detail architecrure shown in Figure\ref{DFblock}\\
\begin{figure}[htbp]
    \centering
    \includegraphics[width=1.0\textwidth]{upblock.png}
    \caption{The architecture of the DFblock: 2 layers linear mlp, tensor from the previous layers}
    \vspace{-15pt}
    \label{DFblock}
\end{figure}
The attention model has two input: the word feature $e\in {R} ^{D \times{N}}$ and the image feature from the hidden layer $h\in {R} ^{\hat{D} \times{N}}$. The word features are transformed  into the semantic space of image features by adding a new perceptron layer, $i.e.$, $e^{\prime}=Ue$, where $\mathrm{U}\in {R} ^{\hat{D} \times{D}}$. And the image feature for the previous hidden layer $h_j\in {R} ^{\hat{D} \times{N}}$, Each column of $h$ is a feature vector of the image.  For $j^{th}$ sub-region, its word-context vector is a dynamic representation of word vector relevant to $h_j$ , calcualted by 
\begin{equation}
c_{j}=\sum_{i=0}^{T-1} \beta_{j, i} e_i^{\prime}, 
\quad where,
\beta_{j,i}=\frac{\exp \left(s_{j,i}^{\prime}\right)}{\sum_{k=0}^{T-1} \exp \left(s_{j, k}^{\prime}\right)}
\end{equation}
${s}'_{j,i}=h^{T}_{j}{e}'_{i}$ represents the weight of the model for the $i^{th}$ word at the $j^{th}$ sub-region of the generated image. Then the word-context matrix for the image set $h$ can be defined by $F^{attn}(e,h)$=$(c_{0},c_{1}...,c_{N-1})\in R^{\hat{D}\times N}$. The output of the attention model is image feature $h$ with representation.\\
The discriminator converts images into features through a series of DownBlock. Then the sentence vector will be replicated and concatenated with image features. The discriminator distinguishing generated images from real samples, the discriminator promotes the synthesize images with higher quality and text-image semantic consistency. And the word feature is used at word level to compute loss by measures the similarity with image-text.\\
The text encoder is a bi-directional Long Short-Term Memory(LSTM). and we directly use the pre-trained model provided by AttnGAN.\\
The image encoder is a Convolutional Neural Network(CNN) that maps images to semantic vectors. we directly use the Inception-v3 model pretrained on ImageNet.
To gnerate realistic images with multiple levels(i.e.,sentence level and word level) of conditions,the loss function of our network is defined as:
\begin{equation}
    \mathcal{L} =\mathcal{L} _{G}+ \gamma \mathcal{L} _{D}
\end{equation}
where
\begin{equation}
\begin{aligned}
        &\mathcal{L} _{G} = -E_{G(z)}\sim\ p_{g}[D(G(z)),e]\\
        &\mathcal{L}_{D} = \lambda \mathcal{L}^{s}+(1-\lambda) \mathcal{L}^{w}.
\end{aligned}
\end{equation}
Here $\gamma$  and $\lambda$ is hyperparameter to balance the terms of equation.
Here we defined those two loss function, respectively.
\begin{equation}
\begin{aligned}
\mathcal{L}^{s}= & - E _{x \sim P _r}[\min (0,-1+D(x, e))] \\
& -(1 / 2) E _{G(z) \sim P _g}[\min (0,-1-D(G(z), e))] \\
& -(1 / 2) E _{x \sim P _{m i s}}[\min (0,-1-D(x, e))] \\
& +k E _{x \sim P _r}\left[\left(\left\|\nabla_x D(x, e)\right\|+\left\|\nabla_e D(x, e)\right\|\right)^p\right] \\
L_G= & - E _{G(z) \sim P _g}[D(G(z), e)] \\
\end{aligned}
\end{equation}
and
\begin{equation}
    \begin{aligned}
        \mathcal{L}^{w} = -\sum^{M}_{i=1} log \frac{exp(\mu \times sim(Q_{i},D_{i}))}{\sum^{M}_{j=1}exp(\mu \times sim(Q_{i},D_{j}))}
        \end{aligned}
\end{equation}
for a batch of image-sentence pairs $\{(Q_{i},D_{i})\}^{M}_{i=1}$, and $\mu$ is a smoothing factor determined by experiments. In this batch of sentences, only $D_{i}$ matches the imaged $Q_{i}$, and treat all other $M-1$ sentences as mismatching descriptions. $sim(Q_{i},D_{i})$ is the cosine similarity between $Q_{i}$ and $D_{i}$ which equal to $\frac{Q_{i}^{T}D_{i}}{\mid \mid Q_{i}\mid \mid \mid \mid D_{i}\mid \mid  } $.

\section{Experiment}

\textbf{Datasets.} We chose two datasets, CUB bird \cite{WahCUB_200_2011} and COCO \cite{lin2014microsoft}, as our primary datasets referring to previous work. The CUB dataset comprises 11,788 bird images across 200 species, each with ten language descriptions. Meanwhile, the COCO dataset has 80,000 images for training and 40,000 for testing, each with five language descriptions.

\noindent \textbf{Training Details.} The optimizer of our SEAttnGAN generator using Adam with $\beta_{1}$=0.0, $\beta_{2}$=0.9 and leaning rate is 0.0001 \cite{kingma2014adam}. The optimizer of our SEAttnGAN discriminator using Adam with $\beta_{1}$=0.0, $\beta_{2}$=0.9 and leaning rate is 0.0004.

\noindent \textbf{Evaluation Details.} We choose the IS (Inception Score) and FID (Fréchet Inception Distance) to evaluate the quality of the generated images to evaluate our work. IS is evaluated by calculating the KL(Kullback–Leibler) divergence scatter of a conditional distribution and a marginal distribution, the higher the IS value then the higher the diversity entropy and the lower the quality entropy then the better the generative model \cite{salimans2016improved}.
\begin{equation}
    IS(G)=exp(E_{x\sim p_{g}}[D_{KL}(p(y\mid x)\mid \mid p(y))])
\end{equation}
FID distance calculates the distance between the real sample, the generative sample in the feature space \cite{heusel2017gans}. The Inception network is first used to extract the features, then the feature space is modelled using a Gaussian model, and then the distance between the two features is solved for. A lower FID means higher quality and diversity of the images.
\begin{equation}
    d^{2}((m_{w},C),(m_{w},C_{w}))=\mid \mid m-m_{w}\mid \mid^{2}_{2}+ Trace(C+C^{W}-2(CC_{W})^{\frac{1}{2}})
\end{equation}

NoP (number of parameters) is used to measure the model size of the tested model, and we take some effective alternatives in the models that are not directly generated.
\vspace{-10pt}
\subsection{Quantitative Evaluation}

In this section, We compare the results generated by StackGAN\cite{zhang2017stackgan}, StackGAN++\cite{zhang2018stackgan++}, MirrorGAN\cite{qiao2019mirrorgan}, SD-GAN\cite{yin2019semantics}, DM-GAN\cite{zhu2019dm}, AttnGAN\cite{xu2018attngan}, CPGAN\cite{liang2020cpgan}, TIME\cite{DBLP:conf/aaai/LiuSZME21}, XMC-GAN\cite{zhang2021cross}, DAE-GAN\cite{ruan2021dae} and SEAttnGAN(ours). We chose CUB bird\cite{WahCUB_200_2011} and COCO datasets\cite{lin2014microsoft} for our IS and FID metric experiments. With the help of the pre-trained model generated by training to represent our NoP, SD-GAN\cite{yin2019semantics}  uses the third-party trained COCO pre-trained model\cite{rombach2021highresolution}, CPGAN\cite{liang2020cpgan}  uses extra pre-trained YOLO v3\cite{farhadi2018yolov3}, XMC-GAN\cite{zhang2021cross}uses extra pre-trained VGG-19\cite{simonyan2014very} and Bert\cite{devlin2018bert},  DAE-GAN uses extra NLTK POS tagging and manually designs rules and TIME uses extra 2-D positional encoding. Where the horizontal line is the experiment designed when it was not proposed by the original paper.
\begin{table}[htbp]
        \caption{The results of FID, IS and NOP comparison on dataset CUB bird and COCO.}
	\begin{center}
	    \begin{tabular}{l|c|c|c|c|c}
		\hline {Model}  & \multicolumn{2}{|c|}{ CUB } & \multicolumn{2}{|c|}{ COCO } & \multicolumn{1}{|c}{Time Consuming}\\
		\cline { 2 - 5 } & IS $\uparrow$ & FID $\downarrow$ & FID $\downarrow$ & NoP $\downarrow$ \\
		\hline StackGAN\cite{zhang2017stackgan}  & 3.70 & - & - & - & - \\
		StackGAN++\cite{zhang2018stackgan++}& 3.84 & - & - & - & -\\
		MirrorGAN\cite{qiao2019mirrorgan}  & 4.56 & 18.34 & 34.71 & - & - \\
		SD-GAN\cite{yin2019semantics}  & 4.67 & - & - & $335 M$\cite{rombach2021highresolution} & 6.18s\\
		DM-GAN\cite{zhu2019dm}  & 4.75 & 16.09 & 32.64 & $46 M$ & 7.06s\\
		\textcolor{blue}{AttnGAN}\cite{xu2018attngan} & \textcolor{blue}{4.36} & \textcolor{blue}{23.98} & \textcolor{blue}{35.49} & \textcolor{blue}{$230M$} & \textcolor{blue}{4.05s}\\
		\hline CPGAN\cite{liang2020cpgan}  & - & - & 55.80 & $318 M$ & -\\
		TIME\cite{DBLP:conf/aaai/LiuSZME21} & 4.91 & 14.30 & 31.14 & $120 M$ & -\\
		XMC-GAN\cite{zhang2021cross} & - & - & 9.30 & $166 M$ & -\\
		DAE-GAN\cite{ruan2021dae}  & 4.42 & 15.19 & 28.12 & $98 M$ & -\\
		\hline \textbf{SEAttnGAN(ours)} & \textbf{4.32} & \textbf{15.03} & \textbf{34.78} & \textbf{26.37M} & \textbf{0.11s} \\
		\hline
	    \end{tabular}
            \vspace{-20pt}
	\end{center}
    \label{tab1}
\end{table}

As shown the results in Table \ref{tab1}, our work has yielded outstanding results. Our motivation is fully completed. From a time complexity perspective, our model takes an order of magnitude less time compared to attn-GAN (4.05s for attn-GAN, and 0.11s for SEAttnGAN). In the test with mainstream generation models in limited data sets, our SEAttnGAN model only uses $26.37M$, which achieves the image generation effect and generation quality of other models more than $ 100M$ or even more than $300M$. Even in the CUB bird's FID item, we beat our improved object AttnGAN with one-tenth of the model size, which is an exciting result. However, our SEAttnGAN did not achieve an absolute advantage in COCO, which may be due to the difficulty of training the generative model in a short time caused by the wide variety of items in the COCO dataset. The longitudinal comparison of the CUB bird dataset of all models is generally better than that of the COCO dataset test.
\vspace{-10pt}
\subsection{Qualitative Evaluation}
\vspace{-24pt}
\begin{figure}[ht]  
    \includegraphics[width=1.0\textwidth]{Q_com.png}
    \vspace{-10pt}
    \caption{The architecture of the SEAttnGAN. FC: fully connected layer. UPBlock: upsample + residual block + DFBlock. DownBlock: downsample + residual block}
      \label{Qual-res}
\vspace{-10pt}
\end{figure}

We compared the visualisation results of our model with three different models in figure \ref{Qual-res}: AttnGAN \cite{xu2018attngan}, DF-GAN, and DM-GAN \cite{zhu2019dm}. The results obtained from AttnGAN and DF-GAN are not as real as those from DM-GAN and our model. The distinction between the bird and the background is less pronounced, particularly in the legs and tail areas. Furthermore, the feathers on the bird's back lack perfect clarity. Moreover, a problem exists in the images generated by DM-GAN; part of the image is missing. In contrast, our model's generated image exhibits a notable characteristic: a high degree of distinction between the bird and the background. The patterns on the bird in our model are easily discernible. Additionally, the background in our model's images has a stronger brightness than those generated by the other three models.

The images of the four columns on the left are all distorted on three models on the COCO dataset, which means that some images are missing, repeated, or folded. However, the background in our model's images showcases a stronger brightness compared to the images generated by the other three models, adding to the visual appeal and overall aesthetics of the output.


\vspace{-10pt}
\section{Conclusion}
\vspace{-10pt}
Based on the results above, it appears that AttnGAN utilizes two attention models during the fine-grained image generation process, which seems unnecessary. However, after implementing a simple and effective target improvement, the model's size was significantly reduced without any significant decrease in performance indicators like FID and IS when compared to the original AttnGAN version. Even in the CUB bird's FID item, we beat our improved object AttnGAN with one-tenth of the model size, which is an exciting result. Our improvement involved eliminating one set of attention models and adding the upsample and residual block, resulting in better image quality. Additionally, our combination of loss functions for the three stages of SEAttnGAN has proven to be effective.
\vspace{-10pt}
\section{Future work}
\vspace{-8pt}
The limitation of our model mainly reflects on the low performance of the COCO dataset. It occasionally generates defective images, meaning it needs to run the model several more times to generate the correct image. Future work will focus on taking methods to overcome these shortcomings. Initially, the number of parameters in our model will be minimised to render our model simple and effective. We anticipate improved performance and more reliable image generation by reducing the complexity. our attention will be directed towards enhancing the Attention-GAN model itself. We will explore and implement novel techniques to refine and optimise its architecture, enhancing its ability to generate high-quality images. Eventually, the Simple and effective philosophy can be extended to most models that are large but widely used by industry.
\vspace{-10pt}
\section{Authorship Statement}
\vspace{-8pt}
All authors contributed to the study's conception and design. Mingyu Jin and Chong Zhang completed the construction and debugging of the basic network framework of SEAttnGAN. Simple and effective improvements to AttnGAN ideas and improvements to the Attention model were completed by Chong Zhang. Loss function reconstruction and SEAttnGAN Mathematical interpretation by Qinkai Yu. Optimising the Attention Model combined with semantic information was done by Haochen Xue. All authors wrote the manuscript, and all authors commented on previous versions.  All authors read and approved the final manuscript.

\vspace{50pt}
\bibliographystyle{plain}
\bibliography{main}

\end{document}
